\documentclass[a4paper]{article}
\input{../../../../../newpreamble.tex}
\title{Math 210A: Homework 4}
\author{Lance Remigio}
\begin{document}
\maketitle
\begin{problem}
   Prove that every cyclic group is abelian. 
\end{problem}

\begin{proof}
Let \( H = \langle x \rangle \) be a cyclic group. Our goal is to show that for all \( g,h \in H \), \( gh = hg \). Indeed, let \( g,h \in H  \). Since \( H  \) is cyclic, we have \( g = x^{a} \) for some \( a \in \Z^{+} \) and \( h = x^{b} \) for some \( b \in \Z^{+} \). Then   
\[  gh = x^{a} \cdot x^{b} = x^{a + b} = x^{b + a} = x^{b} \cdot x^{a} = hg. \]
Hence, \( H  \) is cyclic. If \( H  \) is finite, then let \( | H  |  = n  \) and so we have that the powers \( a  \) and \( b  \) are between \( 0 \) and \( n - 1  \). Employing the same argument above will get our desired result.
\end{proof}

\begin{problem}
    Let \( x \in G  \) where \( G  \) is a group. Prove that \( x^{k} = e  \) if and only if \( | x  |  \mid k \).
\end{problem}
\begin{proof}
\( (\Longrightarrow) \) Suppose \( x^{k} = e  \) and let \( | x  |  = n  \). Using the division algorithm, we can find a unique \( q \in \Z  \) and \( 0 \leq r < n  \) such that \( k = n q + r  \). Then 
\begin{align*}
    e = x^{k} =  x^{qn+r} &= (x^{n})^{q} \cdot x^{r} \tag{\( x^{n} = 1  \)} \\ 
                     &= 1^{q} \cdot x^{r} \\
                     &= x^{r} \tag{\( 0 \leq r < n \)}.
\end{align*}
Note that we must have \( r = 0 \) because otherwise if \( r \neq 0  \), we would have \( x^{r} = e  \) and we would contradict the minimality of \( n  \). Hence, \( | x  |  \mid k \). 

\( (\Longleftarrow) \) Since \( | x  |  \mid k  \), there exists a \( q \in \Z  \) such that \( k = n q  \). Hence, 
\[  x^{k} = x^{n q} = (x^{n})^{q} = 1^{q} = 1. \]
Hence, \( x^{k} = 1  \) which is our desired result.
\end{proof}

\begin{problem}
    Let \( \varphi : G \to H  \) be a homomorphism of groups with kernel \( K  \) and let \( a \in H  \). Prove that if \( u \in \varphi^{-1}(a) \), then \( Ku = \varphi^{-1}(a) \).
\end{problem}
\begin{proof}
Suppose \( u \in \varphi^{-1}(a) \). Then \( a = \varphi(u) \). We will show that \( Ku \subseteq \varphi^{-1}(a) \) and \( Ku \supseteq \varphi^{-1}(a) \). 

Let \( y \in Ku \). Then \( y = k u  \) for some \( k \in K  \). Our goal is to show that \( y \in \varphi^{-1}(a) \). That is, we want to show that \( \varphi(y) = a  \). Indeed, we have 
\[  \varphi(y) = \varphi(ku)  = \varphi(k) \varphi(u) = {1}_{G} a = a. \]
Hence, \( y \in \varphi^{-1}(a) \) and so we have the first containment (\( Ku \subseteq \varphi^{-1}(a)  \)).

For the second containment, let \( y \in \varphi^{-1}(a) \). Then \( \varphi(y) = a  \). Notice that 
\[  y = y {1}_{G} = y(u^{-1} u) = (y u^{-1}) u. \]
Our goal is to show that \( y u^{-1} \in K \). Indeed, since \( u \in \varphi^{-1}(a) \), we have \( \varphi(u) = a  \) implies
\begin{align*}
    \varphi(y u^{-1}) &= \varphi(y) \varphi(u^{-1}) \\
                      &= \varphi(y) (\varphi(u))^{-1} \\
                      &= a a^{-1} \\
                      &= {1}_{H}. 
\end{align*}
Hence, \( y u^{-1} \in K  \) which is our desired result. Thus, \( \varphi^{-1}(a) \subseteq Ku \). 

Since both containments hold, we conclude that \( Ku = \varphi^{-1}(a) \).
\end{proof}

\begin{problem}
    Prove that the subgroup of \( {S}_{4} \) is generated by \( {S}_{4} \) generated by \( (1 \ 2) \) and \( (1 \ 3) (2 \ 4) \) is a noncyclic group of order \( 4  \).   
\end{problem}
\begin{solution}
Observe that none of the elements in the set
\[ \{e, (3  \ 4),  (1 \ 2) , (1 \ 3)(3 \ 4) \}  \]
have orders that are equal to the order of the group above (\( 4 \)). Hence, \( H  \) cannot be a cyclic group. 
\end{solution}

\end{document}

